\input{../tex-inputs/latex-leseansicht-vorspann}

\section[Arthur Schnitzler an Gustav Schwarzkopf, 2{[}6{]}. 6. 1899]{L04422 Arthur Schnitzler an Gustav Schwarzkopf, 2{[}6{]}. 6. 1899}
\nopagebreak\mylabel{L04422v}
\rehead{ }\normalsize\beginnumbering\briefempfaengerindex{Schwarzkopf, Gustav@\textsc{Schwarzkopf, Gustav}!zzzSchnitzler, Arthur@\emph{von Arthur Schnitzler}!1899-06-262@{2{[}6{]}. 6. 1899}|(be}
\toendnotes[C]{\smallbreak\pagebreak[2]}
\correspDesc{Versand  durch Arthur Schnitzler am 2[6]. 6. 1899 in Orahovica
\newline{}Übermittlung  am 26. 6. 1899 in Orahovica
\newline{}Zustellung  am 27. 6. 1899 in Wien
\newline{}Erhalt  durch Gustav Schwarzkopf am 27. 6. 1899 in Wien}\toendnotes[C]{\smallbreak}
\Standort{DLA, A:Schnitzler, HS.1985.1.1897.}
\physDesc{Bildpostkarte, 124 Zeichen
\newline{}Handschrift: Bleistift, deutsche Kurrent
\newline{}Versand: 1) Stempel: »\nobreak{}\oindex{Orahovica@\textbf{Orahovica}|pwk}Orahovica, 26. 6. 99\nobreak{}«.   2) Stempel: »\nobreak{}\oindex{I., Innere Stadt@\textbf{I., Innere Stadt}|pwk}Wien 1/1 1, 27. 6. 1899, 11½V–1N, Bestellt\nobreak{}«. }\toendnotes[C]{\smallbreak}\pstart{}{\pb}Herrn \textsc{Gustav Schwarzkopf}\pend{}\pstart{}Wien\oindex{Wien@\textbf{Wien}|pw}\pend{}\pstart{}\textsc{\textsc{I. Tiefer Graben 23}}\oindex{Wien@\textbf{Wien}!I., Innere Stadt@\textbf{I., Innere Stadt}!Tiefer Graben 23@\textbf{Tiefer Graben 23}|pw}\pend{}{\bigskip}
\pstart
           \centering{}{\pb}\textcolor{gray}{\textbf{\label{K_L04422-1v}\edtext{Pozdrav iz Orahovice\oindex{Orahovica@\textbf{Orahovica}|pw}}{\lemma{\textnormal{\emph{Pozdrav iz Orahovice}}}\Cendnote{\textnormal{kroatisch: Grüße aus
                  Orahovica}}}\label{K_L04422-1}.}}\pend
           \vspace{1em}
\pstart
           \noindent{}{\pb}Herzliche Grüße! Aber nicht der Mond { }ſcheint – ſondern der Regen.\pend
           \pstart Ihr \spacefill\mbox{Arth Sch.}\pend{}\selectlanguage{ngerman}\endnumbering\briefempfaengerindex{Schwarzkopf, Gustav@\textsc{Schwarzkopf, Gustav}!zzzSchnitzler, Arthur@\emph{von Arthur Schnitzler}!1899-06-262@{2{[}6{]}. 6. 1899}|)be}\mylabel{L04422h}
\begin{anhang}
\end{anhang}\newcommand{\dateiname}{L04422}\newcommand{\titel}{Arthur Schnitzler an Gustav Schwarzkopf, 2[6]. 6. 1899}\newcommand{\editorInnen}{Herausgegeben von Jahnke, SelmaMüller, Martin Anton}\input{../tex-inputs/latex-leseansicht-abspann}
